%
% Template for DCS course projects
%
\documentclass[a4paper,11pt,oneside]{book}
\usepackage[latin1]{inputenc}
\usepackage[english]{babel}
\usepackage{amsfonts}
\usepackage{amsmath}
\usepackage{amssymb,amsmath,color}
\usepackage{cite}
\usepackage{graphicx}
\usepackage{float}

\begin{document}
\pagestyle{myheadings}

%%%%%%%%%%% Cover %%%%%%%%%%%
\thispagestyle{empty}                                                 
\begin{center}                                                            
    \vspace{5mm}
    {\LARGE UNIVERSIT\`A DI BOLOGNA} \\                       
      \vspace{5mm}
\end{center}
\begin{center}
  \includegraphics[scale=.27]{figs/logo_unibo}
\end{center}
\begin{center}
      \vspace{5mm}
      {\LARGE School of Engineering} \\
        \vspace{3mm}
      {\Large Master Degree in Automation Engineering} \\
      \vspace{20mm}
      {\LARGE Industrial Robotics} \\
      \vspace{5mm}{\Large\textbf{Laboratory Report}}                  
      \vspace{15mm}
\end{center}
\begin{flushleft}                                                                              
     {\large Professors: \textbf{\@ Claudio Melchiorri}\\  \hspace{20mm} \textbf{\@ Gianluca Palli}} \\  
     
      \vspace{13mm}
\end{flushleft}
\begin{flushright}
      {\large Group 16}\\
\end{flushright}        %capoverso allineato a destra
\begin{center}
\vfill
      {\large Academic year \@2018/2019} \\
\end{center}



\newpage
\thispagestyle{empty}

%%%%%% ABSTRACT %%%%%%%%%%
%\begin{center}
%\chapter*{}
%\thispagestyle{empty}
%{\Huge \textbf{}}\\
%\vspace{15mm}
%\end{center}
%Put here the abstract of the project. 

%%%%%%%%%%%%%%%%%%%%%%%%%%%

\tableofcontents \thispagestyle{empty}
\listoffigures\thispagestyle{empty}



%%%%%%%%% CAPITOLO  %%%%%%%%%%%%%%%%
\chapter{Arnold}

\section*{Introduction}

Arnold is the mobile robot that we have designed and build to perform the Industrial Robotics Laboratory Race. The mechanical structure is based on LEGO parts and uses electric motors and some sensors to move and receive informations whitin the environment. The control is performed by a NXT board. The race basically consists in two parts:
\begin{itemize}
	\item \textit{ \textbf {Line Following:} the goal is to follow a path starting from a square box};
	\item \textit{\textbf {Obstacle Avoidance:} the goal is to move toward an arena containg obstacles};

\end{itemize}
\section{Structure and Hardware}
\subsection {Mechanical assembling}

The mechanical assmebling of the various parts and components of Arnold is the one that we show in the next figure containing different views:\\

\begin{center}
	\includegraphics[scale=0.15]{figs/arnold4}
\end{center}


This structure has been chosen because it guarantees mechanical stability and symmetricity.\\


\subsection {Sensors}

Arnold uses the following sensors:
\begin{itemize}
	\item \textit{ \textbf {Gyroscope:}measures the angular velocity wrt its axis so properly positionating it can be used to determine the robot orientation};
	\item \textit{\textbf {Light Sensor:} measures the light intensity};
	\item \textit{\textbf {Sonar:}determines the distance of eventual obstacles wrt the sensor sending ultrasonic signals and measuring the time in which they come back};
	
\end{itemize}

\section{Algorithms}
\subsection {Threads}

We have used a multi-thread structure.

\subsection {Line Following}

The Line Following part is basically performed by using a PID control.

\subsection {Obstacle Avoidance}

The basic idea of the obstacle avoidance algorithm that we have implemented is the potential gradient method in which we define an attractive potential realated to the robot orientation and a repulsive potential related to the obstacles distances.

\section{Performances}






%%%%%%%%% CAPITOLO  %%%%%%%%%%%%%%%%
\chapter{Matlab Simulations}

\section{Section title}
\subsection {Subsection title}



%%%%%%%%%%%%%%%%%%%%%%%%%%%%%%%%%%%%%



% %%%% APPENDIX %%%%%
% \appendix
% \chapter{Appendix title}
% %%%%%%%%%%%%%%%%%%%%

%%%%%%%%%% BIBLIOGRAPHY %%%%%%%%%%%%%%
\bibliography{bibliography}{}
\bibliographystyle{plain}
\addcontentsline{toc}{chapter}{Bibliography}
%%%%%%%%%%%%%%%%%%%%%%%%%%%%%%%%%%%%%%

\end{document}